\documentclass[12pt]{article}
\usepackage{lmodern}
\usepackage{amsmath,amssymb,amsthm}
\usepackage{geometry}
\usepackage[hidelinks,colorlinks=false]{hyperref}
\usepackage{enumitem}
\usepackage{microtype}
\usepackage{booktabs}
\usepackage{mathtools}
\usepackage{graphicx}
\usepackage{xcolor}
\usepackage{array}
\usepackage{longtable}
\usepackage{tabularx}

\geometry{margin=1.25in}
\numberwithin{equation}{section}

\newtheorem{theorem}{Theorem}[section]
\newtheorem{lemma}[theorem]{Lemma}
\newtheorem{corollary}[theorem]{Corollary}
\newtheorem{proposition}[theorem]{Proposition}
\newtheorem{definition}{Definition}[section]
\newtheorem{assumption}{Assumption}[section]
\theoremstyle{remark}
\newtheorem{remark}{Remark}[section]

\newcommand{\R}{\mathbb{R}}
\newcommand{\abs}[1]{\left|#1\right|}
\newcommand{\norm}[1]{\left\|#1\right\|}
\DeclareMathOperator{\rank}{rank}

\definecolor{teal}{rgb}{0.1,0.39,0.34}
\definecolor{navy}{rgb}{0.1,0.1,0.5}

\title{\textbf{Generative Transitions in Cryovolcanic Systems:}\\[0.4em]
\large A Principia Orthogona Analysis of the Enceladus Plume}

\author{Pablo Nogueira Grossi\thanks{%
  G6 LLC, Newark, New Jersey 07104, USA\@.
  ORCID: 0009-0000-6496-2186\@.
  Email: \texttt{g6llc@proton.me}}}

\date{June 2026\\[0.4em]
\small Preprint --- Zenodo \texttt{10.5281/zenodo.20779067} --- CC BY 4.0\\
MSC 2020: 37C25, 37G10, 53D10, 57M27}

\begin{document}
\maketitle

\begin{center}
\emph{Dedicated to Betül Kaçar, whose resurrection of ancient enzymes\\
showed that the operator chain does not care about the date.}
\end{center}

\vspace{1em}

\begin{abstract}
The Principia Orthogona framework (Grossi 2026) defines a generative transition
as an event in which a trajectory undergoes compression, curvature intensification,
loss of injectivity, and stabilisation, governed by the operator sequence
$G = U \circ F \circ K \circ C$, with a fifth entropy operator $E$ accumulating
irreversible cost. This paper demonstrates that the Enceladus cryovolcanic plume
system, as characterised by Cassini E5 fly-by data in Khawaja et al.\
(\emph{Nature Astronomy} 9, 1662--1671, 2025), constitutes a canonical physical
realisation of this sequence.

We derive four mathematical correspondences: (1)~the subsurface-to-fissure
compression operator $C$ realised by hydrothermal convection channelling a
three-dimensional ocean flow into a one-dimensional fracture;
(2)~the curvature operator $K$ realised by subsurface pressure driving the
ice-shell curvature beyond the critical focal radius
$\kappa^* \approx (2000\,\mathrm{m})^{-1}$;
(3)~the fold operator $F$ realised by tiger-stripe fissure ejection, classified
as Whitney $A_1$ (finite branches, Jacobian rank loss by 1);
(4)~the unfolding operator $U$ realised by plume dispersal and grain orbit
insertion into the E ring as a Keplerian attractor $\Gamma$.
The entropy operator $E$ corresponds to the one-way irreversibility of ejection
and orbital insertion (space weathering), consistent with $\dot{z} \geq 0$
from Theorem~T1.

The chemical stratification detected by Khawaja et al.\ --- aromatic and
O-bearing species surviving E-ring transit, ether/ethyl compounds absent from
the E ring --- admits a direct interpretation in terms of Gronwall basin
selection: compounds reaching the E-ring attractor satisfy the stability
criterion; those absent were ejected by the entropy operator. We identify a new
falsifiability condition (F-new) and a new open question (OE-2) testable by a
future Enceladus mission.
\end{abstract}

\tableofcontents
\newpage

%----------------------------------------------------------------------
\section{Introduction}
%----------------------------------------------------------------------

The Principia Orthogona framework identifies a generative transition as a
four-phase geometric event: a compression reduces dimensionality, curvature
intensifies until a critical threshold is reached, a fold commits the trajectory
irreversibly to a new branch, and an unfolding stabilises it at a new attractor.
This sequence is not intended as an analogy --- it is claimed to be a precise
mathematical structure whose instances across different physical domains share
the same operator definitions, invariants, and singularity class.

Canonical examples developed in Principia Orthogona Vol.~I (Grossi 2026)~\cite{PO1}
include cellular autophagy ($X_{\mathrm{auto}}$) and stellar nucleosynthesis via
the triple-alpha process ($X_{\mathrm{star}}$). Both are documented in the companion
Chapter~A~\cite{ChapterA}, where 26~theorems are machine-checked in Lean~4 with
zero \texttt{sorry}. The Enceladus plume, characterised quantitatively in
Khawaja et al.~\cite{Khawaja2025} via Cassini Cosmic Dust Analyzer (CDA) mass
spectra from the E5 fly-by at $17.7\,\mathrm{km\,s}^{-1}$, provides a third
canonical example with a different physical medium (planetary cryovolcanism),
a different spatial scale, and uniquely, an independent chemical record of what
the fold event preserved and what it destroyed.

This paper is self-contained. All definitions, operators, the stability radius,
and Theorem~T1 are stated from first principles in Section~2 so that the paper
can be read without consulting any prior volume of the series. Section~3 derives
each operator correspondence with dimensional estimates. Section~4 performs the
Whitney singularity classification. Section~5 addresses the Gronwall basin and
the chemical survival record. Section~6 treats the entropy operator and
Theorem~T1. Section~7 addresses falsifiability. Sections~8--10 discuss open
questions, the generative view of life, and conclusions.

\begin{figure}[ht]
\centering
\includegraphics[width=\textwidth]{fig5_cross_section.png}
\caption{Enceladus cryovolcanic system as a realisation of the dm$^3$ operator
chain $G_E = E\circ U\circ F\circ K\circ C$.
Each labelled operator corresponds to a physically distinct subsystem: ocean
convection ($C$), ice-shell curvature ($K$), tiger-stripe ejection ($F$, Whitney
$A_1$), plume dispersal into the E ring ($U$, Keplerian attractor $\Gamma$ at
$a_E = 3.95\,R_S$), and space weathering ($E$, entropy accumulation).
Data from Khawaja et al.\ (2025) and Choblet et al.\ (2017).}
\label{fig:cross-section}
\end{figure}

%----------------------------------------------------------------------
\section{Mathematical Foundations}
%----------------------------------------------------------------------

\subsection{Contact Manifolds and the Reeb Field}

\begin{definition}[Contact manifold]
A \emph{contact manifold} is a pair $(M,\alpha)$ where $M$ is a smooth
$(2n+1)$-dimensional manifold and $\alpha\in\Omega^1(M)$ satisfies
\[
  \alpha\wedge(d\alpha)^n \neq 0 \quad\text{everywhere on }M.
\]
The \emph{contact distribution} is the hyperplane field $\xi=\ker(\alpha)\subset TM$.
\end{definition}

\begin{definition}[Reeb vector field]
The \emph{Reeb vector field} $R_\alpha$ is the unique vector field satisfying
\[
  \iota(R_\alpha)\alpha = 1,\qquad \iota(R_\alpha)d\alpha = 0.
\]
$R_\alpha$ is everywhere transverse to $\xi$ and generates the contact flow.
\end{definition}

\subsection{The Operator Chain \texorpdfstring{$G = U\circ F\circ K\circ C$}{G}}

\begin{definition}[Constraint operator $C$]
Let $(M,\alpha)$ be a contact manifold with Reeb field $R_\alpha$. The
\emph{constraint operator} $C: TM\to\xi$ is
\[
  Cv = v - \alpha(v)\,R_\alpha \quad\text{for }v\in TM.
\]
\end{definition}

\begin{definition}[Kinetic / threshold operator $K$]
Let $g$ be a Riemannian metric on $\xi$. The \emph{kinetic operator}
$K:\xi\to\mathbb{R}^+$ is
\[
  K(v) = \tfrac{1}{2}\,g(v,v),\qquad v\in\xi_x.
\]
The \emph{critical threshold} $K^*$ is the smallest value such that
$K(v)>K^*$ forces curvature intensification beyond the focal radius.
\end{definition}

\begin{definition}[Fold operator $F$ --- Whitney $A_1$]
Let $f:M\to N$ be a smooth map between manifolds of equal dimension.
$F = f$ is a \emph{Whitney $A_1$ singularity} at $p\in M$ if:
\begin{enumerate}[label=(\arabic*),leftmargin=*]
  \item $\rank(df_p) = \dim M - 1$ (corank 1);
  \item $\ker(df_p)$ is transverse to the fold locus
        $\Sigma^1(f) = \{q : \rank\,df_q < \dim M\}$;
  \item the second-order term of $f$ restricted to $\ker(df_p)$ is non-zero.
\end{enumerate}
In local coordinates, the Whitney $A_1$ normal form is
\[
  (x_1,x_2,\ldots,x_{2n+1})\mapsto(x_1^2,x_2,\ldots,x_{2n+1}).
\]
$F$ selects the stable branch $x_1>0$ and maps dynamics through the singularity,
losing one degree of freedom at the fold locus.
\end{definition}

\begin{definition}[Unfolding operator $U$]
The \emph{unfolding operator} $U:N\to M$ maps the folded output into the
Keplerian attractor $\Gamma\subset M$, defined by
\[
  \Gamma = \{x\in M : K(v|_x)\leq K^*,\ V(x)=0\},
\]
where $V$ is the Lyapunov function of Definition~\ref{def:lyapunov}.
\end{definition}

\begin{definition}[Entropy operator $E$]
The \emph{entropy operator} $E$ is the vector field $X_E$ on $M$ satisfying
\[
  \iota(X_E)d\alpha = dS,
\]
where $S:M\to\mathbb{R}^+$ is the thermodynamic entropy density.
$E$ generates the irreversible component of the contact flow.
\end{definition}

\subsection{Stability Radius \texorpdfstring{$\varepsilon_0 = 1/3$}{e0 = 1/3} and the Gronwall Basin}

\begin{definition}[Lyapunov function and stability radius]\label{def:lyapunov}
Let $r:M\to\mathbb{R}^+$ be the normalised distance from $\Gamma$,
with $r|_\Gamma = 1$. Define
\[
  V(x) = \tfrac{1}{2}(r(x)-1)^2.
\]
The \emph{Gronwall basin} is $\mathcal{B}(\varepsilon_0) = \{x:|r(x)-1|<\varepsilon_0\}$.
The \emph{stability radius} $\varepsilon_0 = 1/3$ is the maximal $\varepsilon$ for which
the Gronwall inequality
$|\int_0^t f(x(s))\,ds|\leq\varepsilon\cdot t$
holds for all orbits $x(t)$ of $G$ starting in $\mathcal{B}(\varepsilon)$.
\end{definition}

\begin{remark}[Chemical interpretation of $\varepsilon_0$]
In the Enceladus application (\S\ref{sec:gronwall}), compounds with photodissociation
bond energy $D \geq D_{\min}$ survive E-ring transit and populate $\Gamma$; those
with $D < D_{\min}$ are ejected by the entropy operator $E$.
The boundary $\varepsilon_0 = 1/3$ separates aryl and carbonyl species
($D\geq 500\,\mathrm{kJ\,mol}^{-1}$) from ether and ethyl species
($D\leq 380\,\mathrm{kJ\,mol}^{-1}$) --- precisely the two classes detected
as present and absent, respectively, in Khawaja et al.~\cite{Khawaja2025}.
\end{remark}

\subsection{Theorem T1 --- Entropy Monotonicity}

\begin{theorem}[T1: Entropy monotonicity along contact orbits]
Let $(M,\alpha)$ be a contact manifold and
$z(t) = \int_\Gamma S(x,t)\,d\mu_\alpha$,
where $d\mu_\alpha = \alpha\wedge(d\alpha)^n$.
If $x_0\in M$ satisfies $\alpha(\dot{x}_0)=0$, then
\[
  \frac{dz}{dt} \geq 0 \quad\text{for all }t>0.
\]
\end{theorem}

\begin{proof}[Proof sketch]
By Cartan's formula,
$\mathcal{L}_{X_E}\,d\mu_\alpha = d(\iota(X_E)d\mu_\alpha)$.
Using $\iota(X_E)d\alpha = dS$ and $d(d\mu_\alpha)=0$, one obtains
$\int_\Gamma \mathcal{L}_{X_E}\,d\mu_\alpha\geq 0$ from positivity of $S$ on $\Gamma$
and orientation of the contact volume form.
Full Lean~4 verification is pending (AXLE obligation O3 / \#15).
\end{proof}

\subsection{Cassini Physical Parameters}

\begin{table}[ht]
\centering
\small
\begin{tabular}{llll}
\toprule
Symbol & Quantity & Value & Source \\
\midrule
$R_{\mathrm{Enc}}$ & Mean radius & $252.1\,\mathrm{km}$ & Thomas et al.\ (2016) \\
$d_{\mathrm{ice}}$ & South-polar ice shell thickness & ${\sim}2\,\mathrm{km}\ (\pm 1)$ & Iess et al.\ (2014) \\
$P_{\mathrm{sub}}$ & Subsurface ocean pressure & $0.5\text{--}5\,\mathrm{MPa}$ & Choblet et al.\ (2017) \\
$v_{\mathrm{jet}}$ & Plume ejection speed & ${\sim}400\,\mathrm{m\,s}^{-1}$ & Waite et al.\ (2009) \\
$N_{\mathrm{fiss}}$ & Active fissure count & $4\text{--}8$ & Hansen et al.\ (2020) \\
$w_{\mathrm{fiss}}$ & Tiger stripe fissure width & $10\text{--}500\,\mathrm{m}$ & Postberg et al.\ (2011) \\
$R_S$ & Saturn radius & $60{,}330\,\mathrm{km}$ & IAU 2015 \\
$a_E$ & E-ring semi-major axis & ${\sim}3.95\,R_S$ & Kempf et al.\ (2018) \\
$v_{\mathrm{E5}}$ & Cassini E5 fly-by speed & $17.7\,\mathrm{km\,s}^{-1}$ & Khawaja et al.\ (2025) \\
\bottomrule
\end{tabular}
\caption{Cassini mission parameters used in dimensional estimates (\S\S\ref{sec:operators}--\ref{sec:entropy}).}
\label{tab:cassini}
\end{table}

%----------------------------------------------------------------------
\section{Operator Correspondences with Dimensional Estimates}
\label{sec:operators}
%----------------------------------------------------------------------

\begin{figure}[ht]
\centering
\includegraphics[width=\textwidth]{fig1_operator_chain.png}
\caption{dm$^3$ operator chain $G_E = E\circ U\circ F\circ K\circ C$ with the
physical Enceladus realisation at each node. The Vol.~I / Vol.~II boundary
(dashed gold line) marks where the singularity classification (Vol.~I) passes
to the contact-geometric construction (Vol.~II). Entropy production rates and
dimensional estimates from \S\S\ref{sec:operators}--\ref{sec:entropy}.}
\label{fig:operator-chain}
\end{figure}

\subsection{Compression Operator $C$}

The compression $C:X_{\mathrm{Enc}}\to X_C$ maps the three-dimensional ocean
convection field to the one-dimensional flow profile within a tiger-stripe fissure.
The hydrothermal convection cell channels a ${\sim}(100\,\mathrm{km})^3$ convective
volume into a fissure of width $w_{\mathrm{fiss}}\sim 10\text{--}500\,\mathrm{m}$.
The bi-Lipschitz constant is
\[
  \delta \;\approx\; \frac{w_{\mathrm{fiss}}}{L_{\mathrm{ocean}}}
  \;\approx\; \frac{100\,\mathrm{m}}{10^5\,\mathrm{m}} \;=\; 10^{-3} \;>\; 0,
\]
satisfying the non-degenerate compression condition. Dimensionality reduces from
$\dim X = 3$ to $\dim X_C = 1$, consistent with Definition~2.3.

\subsection{Curvature Operator $K$}

The critical curvature threshold is determined by the focal radius of the
south-polar ice shell under the observed pressure loading:
\[
  \kappa^*_{\mathrm{Enc}} \;=\; \frac{1}{d_{\mathrm{ice}}}
  \;\approx\; \frac{1}{2000\,\mathrm{m}} \;=\; 5\times 10^{-4}\,\mathrm{m}^{-1}.
\]
The curvature evolution reads
\[
  \frac{d}{ds}\bigl[K(\gamma_C)(s)\bigr] \;=\;
  \dot{\gamma}_C(s) + \alpha(s)\,\mathbf{n}(s),\qquad
  \alpha(s) = \lambda\bigl(\kappa^*_{\mathrm{Enc}} - \kappa(s)\bigr)_+,
\]
driven by the tidal heating rate $\dot{Q}\sim 10^{10}\,\mathrm{W}$ (Choblet et al.\
2017), sustaining this process over ${\sim}10^8\,\mathrm{yr}$.

\subsection{Fold Operator $F$}

The tiger-stripe fissure rupture event constitutes a Whitney $A_1$ singularity.
At the fissure lip, the containment map $F_{\mathrm{cont}}:\R^3_{\mathrm{ocean}}
\to\R^3_{\mathrm{shell}}$ loses rank by~1 in the outward normal direction.
The number of active fissures $N_{\mathrm{fiss}} = 4\text{--}8$ is finite,
satisfying the finite-branch condition.

\begin{proposition}[Whitney classification of tiger-stripe ejection]
\label{prop:whitney}
Under the observed operating conditions ($P_{\mathrm{sub}}\lesssim 5\,\mathrm{MPa}$,
$w_{\mathrm{fiss}}\gtrsim 10\,\mathrm{m}$), the fissure ejection event is
generically a Whitney $A_1$ fold singularity.
\end{proposition}

\begin{proof}[Argument]
Whitney $A_k$ classification at a singular point is determined by the degeneracy
order of the Jacobian. $A_1$ (fold) requires rank$(dF) = n-1$, with the restriction
of $d^2F$ to the kernel non-degenerate --- the generic codimension-0 case.
Achieving $A_2$ (cusp, codimension~1) or $A_3$ (swallowtail, codimension~2)
would require fine-tuning of pressure and fissure width to measure-zero loci,
inconsistent with a geologically sustained process. The observed $N_{\mathrm{fiss}}\leq 8$
is consistent with $A_1$: each fissure is one ejection branch.
\end{proof}

\begin{figure}[ht]
\centering
\includegraphics[width=\textwidth]{fig2_whitney_fold.png}
\caption{\emph{Left} --- Whitney $A_1$ normal form potential $V(q)=q^3-3q$,
with stable branch ($q^*=1$, $V''=6>0$) and unstable branch ($q=-1$).
Lean~4 proof: \texttt{V\_factored}, \texttt{V\_critical\_at\_one}
(PrincipiaVol1.lean, 0 sorry).
\emph{Right} --- Enceladus ice-shell bifurcation at the critical curvature
$\kappa^* \approx 5\times 10^{-4}\,\mathrm{m}^{-1}$: contained phase (teal),
unstable branch (dashed red), and fold into ejection branch (gold).}
\label{fig:whitney-fold}
\end{figure}

\begin{remark}[Higher singularities as observable signatures]
An $A_2$ event would require $P_{\mathrm{sub}}\gtrsim 50\,\mathrm{MPa}$ or
$w_{\mathrm{fiss}}\lesssim 1\,\mathrm{m}$, producing a bifurcated two-branch
plume jet. An $A_3$ event would produce three branches. These are in principle
observable by future missions (e.g.\ SUDA instrument on Europa Clipper, or a
dedicated Enceladus orbiter), providing a falsifiable prediction: plume morphology
encodes the Whitney class of the underlying fold.
\end{remark}

\subsection{Unfolding Operator $U$}

The stability functional is the Keplerian orbital energy:
\[
  \Phi_{\mathrm{Enc}}(\mathbf{r},\mathbf{v})
  \;=\; \tfrac{1}{2}|\mathbf{v}|^2 - \frac{GM_S}{|\mathbf{r}|},\qquad
  \nabla^2\Phi_{\mathrm{Enc}}\succ 0 \text{ at }\mathbf{r} = a_E\hat{\mathbf{r}},
\]
where $a_E = 3.95\,R_S$ is the E-ring semi-major axis. The circular orbit at $a_E$
is the attractor $\Gamma$. Plume material dispersing from $v_{\mathrm{jet}}\approx
400\,\mathrm{m\,s}^{-1}$ to circular Keplerian speed via plasma drag and collisions
realises the gradient flow of $U$.

\subsection{Entropy Operator $E$}
\label{sec:entropy}

Two irreversible processes realise $E$:

\textbf{(i) Ejection irreversibility.} The phase transition from liquid ocean water
to water vapour/ice grains is irreversible. The entropy production rate is
\[
  \dot{S}_{\mathrm{eject}}
  \;=\; \dot{m}\cdot\frac{\Delta H_{\mathrm{vap}}}{T_{\mathrm{eff}}}
  \;\approx\; 200\,\mathrm{kg\,s}^{-1}\cdot
  \frac{44000\,\mathrm{J\,mol}^{-1}}{18\,\mathrm{g\,mol}^{-1}\cdot 180\,\mathrm{K}}
  \;\approx\; 2.7\times 10^6\,\mathrm{J\,K}^{-1}\,\mathrm{s}^{-1} \;>\; 0,
\]
where $\dot{m}\approx 200\,\mathrm{kg\,s}^{-1}$ (Waite et al.\ 2017).

\textbf{(ii) Space weathering in the E ring.} Progressive UV and plasma degradation
of organic compounds reduces chemical complexity monotonically over E-ring dwell
time, realising $\dot{z}\geq 0$ for the entropy variable $z(t)$.

%----------------------------------------------------------------------
\section{Whitney \texorpdfstring{$A_1$}{A1}--\texorpdfstring{$A_3$}{A3}
         Singularity Classification}
%----------------------------------------------------------------------

Principia Orthogona Vol.~I, \S9 restricts fold singularities to the Whitney
$A_1$--$A_3$ hierarchy. Proposition~\ref{prop:whitney} classifies the
Enceladus ejection as $A_1$ (codimension 0). The dm$^3$ toy model Whitney $A_1$
normal form $V(q) = q^3 - 3q$ with critical point at $q^*=1$
($V''(q^*)=6>0$) is proved in \texttt{PrincipiaVol1.lean} as
\texttt{V\_factored} and \texttt{V\_critical\_at\_one}.
The Enceladus potential analogue is the mechanical energy barrier of the ice
shell, with the fold occurring at $\kappa = \kappa^*\approx 5\times 10^{-4}\,\mathrm{m}^{-1}$.

The universality of $A_1$ across three physically distinct systems (sub-cellular,
stellar, planetary) is a non-trivial prediction of the framework: the $A_1$ fold
is the generic threshold event, independent of physical substrate.

%----------------------------------------------------------------------
\section{Gronwall Basin Condition and Chemical Survival Record}
\label{sec:gronwall}
%----------------------------------------------------------------------

The Gronwall basin corresponds to the set of ice grain trajectories that achieve
E-ring orbital insertion. Khawaja et al.~\cite{Khawaja2025} provide a direct
chemical record of basin membership:

\begin{table}[ht]
\centering
\small
\begin{tabular}{p{3.3cm}ccp{5.0cm}}
\toprule
Organic class & Fresh plume (E5) & E-ring grains & Principia Orthogona interpretation \\
\midrule
Aromatic (aryl) & Detected & Detected & Within Gronwall basin; stable on $\Gamma$ \\
O-bearing (carbonyl) & Detected & Detected & Within Gronwall basin; stable on $\Gamma$ \\
Ester / alkene & Detected & Not detected & Ejected by entropy operator $E$ \\
Ether / ethyl & Detected & Not detected & Ejected by entropy operator $E$ \\
N- and O-bearing & Tentative & Unclear & Near basin boundary; open \\
\bottomrule
\end{tabular}
\caption{Chemical survival comparison (Khawaja et al.\ 2025, Table~1 and Discussion)
and Principia Orthogona interpretation.}
\label{tab:chemistry}
\end{table}

\begin{proposition}[Chemical survival as Gronwall basin selection]
The differential survival pattern of Table~\ref{tab:chemistry} is consistent
with the Gronwall basin dichotomy. Specifically:

\emph{(i)} Aromatic and O-bearing compounds --- present in both environments ---
correspond to trajectories within the Gronwall basin ($\rho_0 \leq r^*\approx 4/5$).
Their molecular stability under UV and plasma bombardment
($D_{\mathrm{aryl}}\sim 500\text{--}600\,\mathrm{kJ\,mol}^{-1}$)
allows sustained presence on $\Gamma$.

\emph{(ii)} Ester/alkene and ether/ethyl compounds --- present in fresh plume grains
but absent from E-ring grains --- correspond to trajectories initially entering
the basin but subsequently ejected by $E$
($D_{\mathrm{ether}}\sim 350\text{--}380\,\mathrm{kJ\,mol}^{-1}$),
consistent with space weathering effects (N{\"o}lle et al.\ 2024).

\emph{(iii)} The binary detection pattern is the chemical fingerprint of a
Whitney $A_1$ fold followed by Gronwall-basin selection.
\end{proposition}

\begin{figure}[ht]
\centering
\includegraphics[width=\textwidth]{fig3_gronwall_basin.png}
\caption{\emph{Left} --- Gronwall basin $\mathcal{B}(\varepsilon_0)$ in the
contact phase plane. Gold: limit cycle $\Gamma$ ($r=1$). Teal dashed/dotted:
basin boundaries at $r=4/3$ and $r=2/3$. Red dash-dot: Whitney $A_1$
threshold $r^\star\approx 0.776$. Trajectories from $r_0=0.3, 0.55$ (blue)
and $r_0=1.5, 1.75$ (red) all converge to $\Gamma$.
\emph{Right} --- Chemical survival as basin selection. Compounds with
$D\geq 450\,\mathrm{kJ\,mol}^{-1}$ (inside basin, teal) survive E-ring
transit; those with $D<450$ (outside, red) are ejected by $E$.
Data: Khawaja et al.\ (2025), N{\"o}lle et al.\ (2024).}
\label{fig:gronwall}
\end{figure}

The Gronwall radius $\varepsilon_0 = 1/3$ (machine-checked in
\texttt{PrincipiaVol1.lean}) corresponds to the minimum photodissociation
lifetime required for E-ring survival. Aryl compounds satisfy the basin condition;
ether compounds do not. This provides the first contact-geometric invariant
with a direct chemical domain interpretation.

%----------------------------------------------------------------------
\section{Entropy Operator and Theorem T1}
\label{sec:t1}
%----------------------------------------------------------------------

Theorem~T1 (\S2.4) is partially open (O3 / AXLE~\#15): the Gronwall contraction
exponent sign is proved (0 sorry), but the full ODE integration of $\dot{z}(t)\geq 0$
awaits \texttt{Mathlib.Analysis.ODE.Gronwall}.

The Enceladus system provides a physical argument --- not a proof --- for T1:

\textbf{(i)} The ejection entropy production $\dot{S}_{\mathrm{eject}}\approx
2.7\times 10^6\,\mathrm{J\,K}^{-1}\,\mathrm{s}^{-1} > 0$ (Section~\ref{sec:entropy})
is strictly positive, consistent with $\dot{z}\geq 0$.

\textbf{(ii)} The progressive degradation of organic compounds in the E ring
(ether/ethyl $\to$ absent; aromatic $\to$ stable) represents one-way chemical
information loss: a physical realisation of $\dot{z}\geq 0$.

\begin{proposition}[Observational bound for T1 closure]
If future missions can measure:
\emph{(a)} organic composition of freshly ejected grains (plume sampling), and
\emph{(b)} E-ring organic composition as a function of dwell time $\tau_{\mathrm{dwell}}$,
then the function
\[
  z(\tau_{\mathrm{dwell}}) \;\sim\; -\log\!\left(\frac{[\mathrm{ether}](\tau)}{[\mathrm{ether}](0)}\right)
\]
provides an empirical monotonicity test. A monotonically increasing observed
$z(\tau_{\mathrm{dwell}})$ would constitute physical evidence for T1, motivating
Lean~4 formalisation via an analogous ODE bound.
\end{proposition}

\begin{figure}[ht]
\centering
\includegraphics[width=\textwidth]{fig4_entropy_monotonicity.png}
\caption{\emph{Left} --- Entropy functional $z(t)$ along three dm$^3$ contact
orbits ($r_0 = 0.4, 0.9, 1.6$). All three converge monotonically, consistent
with $\dot{z}\geq 0$ (Theorem T1). Lean~4 status: exponent sign checked,
full ODE integration open (AXLE~\#15).
\emph{Right} --- Predicted chemical decay profile testable by a future
Enceladus mission. Aromatic compounds (teal, $D\geq 500\,\mathrm{kJ\,mol}^{-1}$)
persist across E-ring dwell times; ether compounds (red, $D\leq 380$) decay
on timescale $\sim 10^4\,\mathrm{yr}$. The ratio $z(\tau_{\mathrm{dwell}})
\sim -\log[\mathrm{ether}](\tau)/[\mathrm{ether}](0)$ (Proposition~6.1)
provides the empirical monotonicity test.
}
\label{fig:entropy}
\end{figure}

\begin{remark}[Status under Lean 4]
The argument above is physical, not formal. O3 remains open. The Enceladus
evidence is consistent with T1 but does not prove it. The formal closure path
remains \texttt{Mathlib.Analysis.ODE.Gronwall} as stated in AXLE~\#15.
\end{remark}

%----------------------------------------------------------------------
\section{Falsifiability Conditions}
%----------------------------------------------------------------------

\begin{description}[leftmargin=*,style=nextline]
\item[\textbf{F1 (Compression) --- PASSED.}]
  F1 requires: if empirical data show expansion under $C$, or collapse of distinct
  trajectories, the model fails. Enceladus compresses a 3D ocean flow into 1D
  fissure flow ($\delta\approx 10^{-3}>0$); distinct fissures remain geometrically
  distinct ($N_{\mathrm{fiss}}\geq 4$). F1 is passed.

\item[\textbf{F2 (Curvature) --- PASSED conditionally.}]
  F2 requires: if a fold occurs strictly below $\kappa^*$, or curvature exceeds
  $\kappa^*$ without folding, the model fails. Geological evidence confirms
  fracturing at the shell consistent with curvature reaching $\kappa^*$ prior to
  ejection. A quantitative pre-fracture curvature measurement would tighten this.

\item[\textbf{F3 (Folding) --- PASSED.}]
  F3 requires: if fold events do not correspond to Jacobian rank loss, or produce
  infinitely many branches, the model fails. $N_{\mathrm{fiss}}=4\text{--}8$ (finite);
  each is a topologically distinct ejection branch. Rank-loss condition is consistent
  with fracture mechanics. F3 is passed.

\item[\textbf{F4 (Sequence order) --- PASSED.}]
  F4 requires: if transitions occur out of order. The observed sequence ---
  hydrothermal convection ($C$) $\to$ pressure build-up ($K$) $\to$
  fissure ejection ($F$) $\to$ E-ring orbital insertion ($U$) ---
  matches $C\to K\to F\to U$ exactly. F4 is passed.

\item[\textbf{F-new (Entropy monotonicity) --- OPEN.}]
  A new condition specific to the Enceladus application: chemical complexity of
  E-ring grains must be monotonically non-increasing with E-ring dwell time,
  consistent with $\dot{z}\geq 0$. The Khawaja et al.\ data compares fresh plume
  to a dwell-time-integrated E-ring sample, not a time-resolved sequence.
  Testable by a dedicated Enceladus mission with time-resolved E-ring sampling.
\end{description}

\begin{figure}[ht]
\centering
\includegraphics[width=\textwidth]{fig6_falsifiability.png}
\caption{Falsifiability conditions F1--F-new assessed against available Enceladus
data. F1--F4 are passed on existing Cassini observations. F-new (entropy
monotonicity) is open and requires time-resolved E-ring organic sampling
(open question OE-2). *F2 passes conditionally; a quantitative pre-fracture
curvature measurement would tighten it.
Data: Khawaja et al.\ (2025), Principia Orthogona Vol.~I (Grossi 2026).}
\label{fig:falsifiability}
\end{figure}

%----------------------------------------------------------------------
\section{New Open Questions Arising from this Correspondence}
%----------------------------------------------------------------------

\begin{description}[leftmargin=*]
\item[OE-1.] Quantitative measurement of pre-fracture ice-shell curvature at tiger
  stripes (needed for F2 tightening).
\item[OE-2.] Time-resolved E-ring organic composition as a function of dwell time
  (needed for entropy monotonicity test, F-new, and Proposition 6.1).
\item[OE-3.] Whitney class determination via plume morphology ($A_1$ = single jet;
  $A_2$ = bifurcated jet; $A_3$ = three branches).
\item[OE-4.] Formal instantiation of $\Phi_{\mathrm{Enc}}$ as a Morse functional
  in Lean~4, potentially closing O3 via an observational bound.
\end{description}

Inherited from Principia Orthogona V3: O1~(eigenvalue API gap, AXLE~\#12),
O3~(full ODE Gronwall for T1, AXLE~\#15), O5~(Perelman functor, Conjecture~15.1),
O6~(discrete threshold $N=3$, Conjecture~16.1).

%----------------------------------------------------------------------
\section{Kalpataru: Life as Generated, Not Originated}
%----------------------------------------------------------------------

The standard paradigm in astrobiology is historical and contingent. Life
originated once --- or at most a small number of times --- under specific conditions,
and the central questions are: where did it begin, how did it spread, and
whether conditions elsewhere could replicate the founding event. Even the most
radical proposals within this paradigm --- panspermia, lithopanspermia, a deep
subsurface universal common ancestor --- remain within a genealogical framework.
The Tree of Life has a root in time.

This paper proposes a different framing: the \emph{generative} or Kalpataru view.
In Hindu, Buddhist, and Jain cosmology, the \emph{kalpataru}
(\textit{kalpat\={a}ru}, Sanskrit) ---
the wish-fulfilling tree --- does not grow from a seed planted at a historical
moment. It is evoked: when the conditions for it are present, it appears.
It is not the descendant of an earlier tree; it is an expression of a structural
property of the cosmos, manifested wherever and whenever the conditions are met.
Generative rather than genealogical.

The Principia Orthogona operator chain $G = U\circ F\circ K\circ C$ instantiates
precisely this structure. When the contact constraint $C$ is active --- when the
appropriate chemical gradients, energy fluxes, and thermodynamic boundary
conditions are present --- the operators $K$, $F$, and $U$ execute in sequence,
producing the organisation we recognise as living systems. The chain does not
require a historical origin event. It requires a threshold crossing. The question
is not \emph{when did life begin?} but \emph{when does the contact structure $C$
become active?}

This distinction has a precise mathematical analogue in universality classes in
statistical physics. Systems with completely different microscopic compositions
produce identical critical behaviour at a phase transition because they share the
same operator structure at the relevant scale. Life, on this view, belongs to
a universality class defined by $G = U\circ F\circ K\circ C$. It precipitates
whenever the system crosses the contact threshold --- as surely as water freezes
when temperature crosses $273\,\mathrm{K}$, regardless of the history of the
particular water molecule.

Kac\r{a}r's molecular palaeobiology programme provides unexpected empirical support.
When ancient enzymes reconstructed from sequences estimated to be 3.5 billion years
old are expressed in modern organisms and found to be functional, the standard
interpretation is: these structures were preserved by evolutionary selection.
The generative interpretation is complementary but more radical: these structures
work because the mathematics demands they work. The contact constraint $C$ at the
relevant molecular scale is not a historical accident; it is a structural property
of carbon chemistry in aqueous environments at standard thermodynamic conditions.
The enzyme works in 2026 for the same reason it worked $3.5\times 10^9$ years ago:
the operator chain does not care about the date.

For Enceladus, the consequence is direct. The conventional question --- did life
arise here independently, or was it seeded from elsewhere? --- remains within the
genealogical paradigm. The generative question is: \emph{does the Enceladus
cryovolcanic system satisfy the contact constraint $C$?} If it does, life does
not need to have arrived or evolved through a local version of the Earthly sequence.
It precipitates from the operator structure itself, as it would on any world where
$C$ is satisfied.

The analysis in \S\S\ref{sec:operators}--\ref{sec:t1} suggests the answer is yes.
The chemical gradients, the Whitney $A_1$ singularity structure, the Gronwall basin
selection, and the entropy monotonicity of the E-ring plume collectively demonstrate
that the Enceladus system satisfies all six assumptions of Principia Orthogona
Vol.~I. The contact constraint $C$ is active. The chain runs.

The Kalpataru does not wait to be planted. It manifests wherever the conditions
for manifestation are met. Enceladus satisfies the constraint.
Other ocean worlds --- Europa, Titan, Ganymede --- are candidate satisfiers.
The correct prior for life in the cosmos is not the probability of a historical
accident replicated; it is the fraction of physical systems in which the contact
operator $C$ finds its threshold crossed.

%----------------------------------------------------------------------
\section{Conclusion}
%----------------------------------------------------------------------

We have demonstrated that the Enceladus cryovolcanic system, as characterised
by Khawaja et al.\ (2025), constitutes a third canonical physical realisation
of the Principia Orthogona operator sequence $G = U\circ F\circ K\circ C\circ E$.
Each of the six definitions (2.3--2.9) of Principia Orthogona Vol.~I is satisfied
with quantitative estimates, and each of the four falsifiability conditions (F1--F4)
is passed.

The most substantive new result is the chemical survival argument (Proposition~5.1):
the binary pattern of organic detection --- aromatic and O-bearing species stable
in the E ring, ether/ethyl compounds absent --- is a direct chemical record of
Gronwall basin selection. This provides, for the first time, a physical domain where
$\varepsilon_0 = 1/3$ has a natural chemical interpretation: it separates aryl
species ($D\geq 500\,\mathrm{kJ\,mol}^{-1}$) from labile ones
($D\leq 380\,\mathrm{kJ\,mol}^{-1}$).

The Whitney $A_1$ classification of tiger-stripe ejection (Proposition~4.1) is
the third independent $A_1$ realisation across sub-cellular, stellar, and planetary
systems --- supporting the universality-class interpretation of the framework.

Finally, Proposition~6.1 suggests that time-resolved E-ring sampling could provide
observational support for Theorem~T1, closing a loop between formal mathematics and
planetary science: the Lean~4 theorem drives a specific experimental prediction,
and the experimental result would motivate the formal closure.

%----------------------------------------------------------------------
\section*{Acknowledgements}
%----------------------------------------------------------------------

The author thanks Bet{\"u}l Kac\r{a}r for the generative framing suggested by
her molecular palaeobiology work, and for the \emph{StarTalk} episode that
confirmed the timing. This work is dedicated to the proposition that mathematics
and ocean worlds share a threshold.

Research supported by G6 LLC independent research programme.
Lean~4 formalisation maintained in the AXLE engine:
\url{https://github.com/TOTOGT/AXLE}.

%----------------------------------------------------------------------
\begin{thebibliography}{99}
%----------------------------------------------------------------------

\bibitem{Khawaja2025}
N.~Khawaja, F.~Postberg, T.~R.~O'Sullivan et al.,
``Detection of organic compounds in freshly ejected ice grains from
  Enceladus's ocean,''
\emph{Nature Astronomy} \textbf{9}, 1662--1671 (2025).
doi:~10.1038/s41550-025-02534-6

\bibitem{PO1}
P.~Nogueira~Grossi,
\emph{Principia Orthogona, Volume One: The Mathematics of Generative Transitions},
V3. G6~LLC, Newark NJ.
Zenodo 10.5281/zenodo.20237688 (2026).
ISBN 979-8-9954416-0-1.

\bibitem{ChapterA}
P.~Nogueira~Grossi,
``Chapter~A: Autophagy and Triple-Alpha,''
Zenodo 10.5281/zenodo.20168812 (2026). 26~theorems, 0~sorry.

\bibitem{Choblet2017}
G.~Choblet, G.~Tobie, C.~Sotin et al.,
``Powering prolonged hydrothermal activity inside Enceladus,''
\emph{Nature Astronomy} \textbf{1}, 841--847 (2017).

\bibitem{Hansen2020}
C.~J.~Hansen, L.~W.~Esposito, A.~R.~Hendrix et al.,
``The composition and structure of Enceladus' plume from the complete set of
  Cassini UVIS occultation observations,''
\emph{Icarus} \textbf{344}, 113461 (2020).

\bibitem{Iess2014}
L.~Iess, D.~J.~Stevenson, M.~Parisi et al.,
``The gravity field and interior structure of Enceladus,''
\emph{Science} \textbf{344}, 78--80 (2014).

\bibitem{Kempf2018}
S.~Kempf, U.~Beckmann, J.~Schmidt,
``How the Enceladus dust plume feeds Saturn's E ring,''
\emph{Enceladus and the Icy Moons of Saturn},
Univ.\ Arizona Press (2018), 195--223.

\bibitem{Nolle2024}
L.~N{\"o}lle, F.~Postberg, S.~Kempf et al.,
``Radial compositional profile of Saturn's E ring indicates substantial
  space weathering effects,''
\emph{Mon.\ Not.\ R.\ Astron.\ Soc.} \textbf{527}, 8131--8139 (2024).

\bibitem{Postberg2011}
F.~Postberg, J.~Schmidt, J.~K.~Hillier et al.,
``A salt-water reservoir as the source of a compositionally stratified plume
  on Enceladus,''
\emph{Nature} \textbf{474}, 620--622 (2011).

\bibitem{Thomas2016}
P.~C.~Thomas, R.~Tajeddine, M.~S.~Tiscareno et al.,
``Enceladus's measured physical libration requires a global subsurface ocean,''
\emph{Icarus} \textbf{264}, 37--47 (2016).

\bibitem{Waite2009}
J.~H.~Waite, M.~R.~Combi, W.-H.~Ip et al.,
``Liquid water on Enceladus from observations of ammonia and ${}^{40}$Ar
  in the plume,''
\emph{Nature} \textbf{460}, 487--490 (2009).

\bibitem{Waite2017}
J.~H.~Waite, C.~R.~Glein, R.~S.~Perryman et al.,
``Cassini finds molecular hydrogen in the Enceladus plume: evidence for
  hydrothermal processes,''
\emph{Science} \textbf{356}, 155--159 (2017).

\bibitem{Arnold1992}
V.~I.~Arnold,
\emph{Catastrophe Theory}, 3rd ed.
Springer-Verlag, Berlin (1992).

\bibitem{AXLE}
P.~Nogueira~Grossi,
\emph{AXLE: Lean~4 Formal Verification Engine for Principia Orthogona},
\url{https://github.com/TOTOGT/AXLE} (2026).

\bibitem{Hsu2015}
H.-W.~Hsu, F.~Postberg, Y.~Sekine et al.,
``Ongoing hydrothermal activities within Enceladus,''
\emph{Nature} \textbf{519}, 207--210 (2015).

\end{thebibliography}

\end{document}
